% Author: Xu Chen, University of Washington
% Lab 04 scene: a phone hung on a string from a door frame, swinging in the plane of its screen, phyphox
% recording the gyroscope; the measurement chain from the swing to zeta and omega_n.
% scene_hw.tex defines \cipHW and inputs this file: same drawing, homework numbers only.
\documentclass[border=4pt]{standalone}
% Author: Xu Chen, University of Washington
% Shared TikZ style for the Control in Practice scene figures. % Author: Xu Chen, University of Washington
% Shared TikZ style for the Control in Practice scene figures. % Author: Xu Chen, University of Washington
% Shared TikZ style for the Control in Practice scene figures. \input{../../../docs/figures/cipstyle} from a lab.
\usepackage{amsmath,amssymb}
\usepackage[T1]{fontenc}
\usepackage{lmodern}
\usepackage{tikz}
\usepackage[siunitx]{circuitikz}
\usetikzlibrary{arrows.meta,calc,positioning,decorations.pathmorphing,decorations.markings,patterns,shapes.geometric,shadings,fit,backgrounds}
% restrained palette
\definecolor{ink}{HTML}{1A1A1A}
\definecolor{steel}{HTML}{8C96A0}
\definecolor{metal}{HTML}{D5DAE0}
\definecolor{accent}{HTML}{1F4E79}   % controller, signals we choose
\definecolor{plantc}{HTML}{FBE7D6}   % plant fill
\definecolor{sensorc}{HTML}{DCE8F3}  % sensor fill
\definecolor{dist}{HTML}{B03A2E}     % disturbances, forces from the environment
\definecolor{act}{HTML}{2E7D4F}      % actuation
\definecolor{note}{HTML}{555555}
\tikzset{
  every picture/.style={line width=0.8pt, >={Stealth[length=2.4mm]}, font=\small},
  thin line/.style={line width=0.5pt},
  thick line/.style={line width=1.2pt},
  dim/.style={thin line, note, {Stealth[length=1.6mm]}-{Stealth[length=1.6mm]}},
  ground/.style={pattern=north east lines, pattern color=steel, draw=none},
  block/.style={draw=ink, rounded corners=1.5pt, minimum height=9mm, minimum width=16mm, align=center, fill=white},
  plant/.style={block, fill=plantc},
  sensor/.style={block, fill=sensorc},
  ctrl/.style={block, draw=accent, text=accent, thick line},
  sum/.style={draw=ink, circle, minimum size=5mm, inner sep=0pt},
  sig/.style={->, ink},
  distsig/.style={->, dist, thick line},
  force/.style={-{Stealth[length=3mm]}, thick line},
  panel title/.style={font=\normalsize\bfseries, anchor=west, ink},
  caption/.style={font=\footnotesize, note, align=left},
}
 from a lab.
\usepackage{amsmath,amssymb}
\usepackage[T1]{fontenc}
\usepackage{lmodern}
\usepackage{tikz}
\usepackage[siunitx]{circuitikz}
\usetikzlibrary{arrows.meta,calc,positioning,decorations.pathmorphing,decorations.markings,patterns,shapes.geometric,shadings,fit,backgrounds}
% restrained palette
\definecolor{ink}{HTML}{1A1A1A}
\definecolor{steel}{HTML}{8C96A0}
\definecolor{metal}{HTML}{D5DAE0}
\definecolor{accent}{HTML}{1F4E79}   % controller, signals we choose
\definecolor{plantc}{HTML}{FBE7D6}   % plant fill
\definecolor{sensorc}{HTML}{DCE8F3}  % sensor fill
\definecolor{dist}{HTML}{B03A2E}     % disturbances, forces from the environment
\definecolor{act}{HTML}{2E7D4F}      % actuation
\definecolor{note}{HTML}{555555}
\tikzset{
  every picture/.style={line width=0.8pt, >={Stealth[length=2.4mm]}, font=\small},
  thin line/.style={line width=0.5pt},
  thick line/.style={line width=1.2pt},
  dim/.style={thin line, note, {Stealth[length=1.6mm]}-{Stealth[length=1.6mm]}},
  ground/.style={pattern=north east lines, pattern color=steel, draw=none},
  block/.style={draw=ink, rounded corners=1.5pt, minimum height=9mm, minimum width=16mm, align=center, fill=white},
  plant/.style={block, fill=plantc},
  sensor/.style={block, fill=sensorc},
  ctrl/.style={block, draw=accent, text=accent, thick line},
  sum/.style={draw=ink, circle, minimum size=5mm, inner sep=0pt},
  sig/.style={->, ink},
  distsig/.style={->, dist, thick line},
  force/.style={-{Stealth[length=3mm]}, thick line},
  panel title/.style={font=\normalsize\bfseries, anchor=west, ink},
  caption/.style={font=\footnotesize, note, align=left},
}
 from a lab.
\usepackage{amsmath,amssymb}
\usepackage[T1]{fontenc}
\usepackage{lmodern}
\usepackage{tikz}
\usepackage[siunitx]{circuitikz}
\usetikzlibrary{arrows.meta,calc,positioning,decorations.pathmorphing,decorations.markings,patterns,shapes.geometric,shadings,fit,backgrounds}
% restrained palette
\definecolor{ink}{HTML}{1A1A1A}
\definecolor{steel}{HTML}{8C96A0}
\definecolor{metal}{HTML}{D5DAE0}
\definecolor{accent}{HTML}{1F4E79}   % controller, signals we choose
\definecolor{plantc}{HTML}{FBE7D6}   % plant fill
\definecolor{sensorc}{HTML}{DCE8F3}  % sensor fill
\definecolor{dist}{HTML}{B03A2E}     % disturbances, forces from the environment
\definecolor{act}{HTML}{2E7D4F}      % actuation
\definecolor{note}{HTML}{555555}
\tikzset{
  every picture/.style={line width=0.8pt, >={Stealth[length=2.4mm]}, font=\small},
  thin line/.style={line width=0.5pt},
  thick line/.style={line width=1.2pt},
  dim/.style={thin line, note, {Stealth[length=1.6mm]}-{Stealth[length=1.6mm]}},
  ground/.style={pattern=north east lines, pattern color=steel, draw=none},
  block/.style={draw=ink, rounded corners=1.5pt, minimum height=9mm, minimum width=16mm, align=center, fill=white},
  plant/.style={block, fill=plantc},
  sensor/.style={block, fill=sensorc},
  ctrl/.style={block, draw=accent, text=accent, thick line},
  sum/.style={draw=ink, circle, minimum size=5mm, inner sep=0pt},
  sig/.style={->, ink},
  distsig/.style={->, dist, thick line},
  force/.style={-{Stealth[length=3mm]}, thick line},
  panel title/.style={font=\normalsize\bfseries, anchor=west, ink},
  caption/.style={font=\footnotesize, note, align=left},
}

\ifdefined\cipHW
  \newcommand{\titletxt}{your phone (homework)}
  \newcommand{\dtxt}{$d$: tape, ${\pm}0.5$\,cm}
  \newcommand{\phonetxt}{$h$, $w$: measure the phone}
  \newcommand{\reltxt}{pull aside 20--30$^\circ$, hold 3\,s,\\ release; record 3\,min}
  \newcommand{\datatxt}{data: \texttt{04\_pendulum\_my\_recording}\\ {\color{note}\scriptsize(no phone only: \texttt{04\_pendulum\_hw\_backup}, simulated, $d=1.05$\,m)}}
\else
  \newcommand{\titletxt}{simulated sample (demos)}
  \newcommand{\dtxt}{$d=0.75$\,m}
  \newcommand{\phonetxt}{$h\times w=0.150\times 0.072$\,m}
  \newcommand{\reltxt}{released from $35^\circ$ after\\ 3\,s still; 3\,min at 100\,Hz}
  \newcommand{\datatxt}{data: \texttt{04\_pendulum\_sample} (simulated)}
\fi
\begin{document}
\begin{tikzpicture}
% ======================= (a) the scene =======================
\node[panel title] at (0,7.75) {(a) The scene: \titletxt};
% door frame
\fill[ground] (0.05,7.0) rectangle (5.05,7.25);
\draw[thick line] (0.05,7.0) -- (5.05,7.0);
\draw[thick line] (0.25,7.0) -- (0.25,1.0);
\draw[thick line] (4.85,7.0) -- (4.85,1.0);
\draw[thin line, steel] (0.05,1.0) -- (5.05,1.0);
\node[font=\footnotesize, note, anchor=north, align=center] at (2.55,0.95) {door frame (or rod, hook);\\ swing over a bed or sofa};
% pivot (hook)
\draw[thick line] (2.2,7.0) -- (2.2,6.88) arc[start angle=90, end angle=-200, radius=0.08];
\coordinate (O) at (2.2,6.8);
\fill (O) circle (0.035);
\node[font=\footnotesize, anchor=east] at (2.05,6.8) {pivot};
% vertical reference and angle
\def\th{22}
\def\dd{3.9}
\draw[thin line, note, dashed] (O) -- ++(0,-5.0);
\coordinate (C) at ($(O)+(-90+\th:\dd)$);
\draw[accent, thin line, ->] ($(O)+(0,-1.35)$) arc[start angle=-90, end angle=-90+\th, radius=1.35];
\node[accent] at ($(O)+(-90+\th/2:1.62)$) {$\theta$};
% string to the top of the phone
\draw[ink] (O) -- ($(C)+(90+\th:0.68)$);
% phone (drawn larger than to scale), rotated by theta about its centre
\begin{scope}[shift={(C)}, rotate=\th]
  \draw[fill=metal, rounded corners=2pt] (-0.34,-0.68) rectangle (0.34,0.68);
  \draw[fill=white, thin line] (-0.27,-0.56) rectangle (0.27,0.56);
  \draw[->, accent, thick line] (0,0) -- (0.62,0) node[right, font=\footnotesize] {$x$};
  \draw[->, accent, thick line] (0,0) -- (0,0.4);
  \node[accent, font=\footnotesize] at (-0.15,0.3) {$y$};
  \draw[accent, thin line, fill=white] (0,0) circle (0.09); \fill[accent] (0,0) circle (0.035);
  \draw[dim] (-0.5,-0.68) -- node[left, font=\footnotesize, ink] {$h$} (-0.5,0.68);
  \draw[dim] (-0.34,-0.84) -- node[below, font=\footnotesize, ink] {$w$} (0.34,-0.84);
\end{scope}
\node[accent, font=\footnotesize] at ($(C)+(-0.1,-0.22)$) {$z$};
% d dimension parallel to the string
\draw[dim] ($(O)+(\th:0.4)$) -- node[sloped, above, font=\footnotesize, ink] {$d$} ($(C)+(\th:0.4)$);
% gravity at the centre and the swing
\draw[force, dist] ($(C)+(0,-0.12)$) -- ++(0,-0.95) node[left, font=\footnotesize] {$mg$};
\draw[<->, thin line] ($(O)+(-90-16:\dd+1.6)$) arc[start angle=-90-16, end angle=-90+26, radius=\dd+1.6];
\node[font=\footnotesize, note, anchor=south] at ($(O)+(-90-16:\dd+1.6)+(0.3,0.08)$) {swing};

% right column: numbers, axes, gyro
\begin{scope}[shift={(5.3,0)}]
\node[font=\footnotesize, align=left, anchor=north west] at (0,7.1)
  {\dtxt\\ (pivot to phone centre)\\ \phonetxt\\[3pt]
   $L_\text{eq}=d+\dfrac{h^2+w^2}{12\,d}$\\[3pt] $\omega_n\approx\sqrt{g/L_\text{eq}}$};
\node[font=\footnotesize, align=left, anchor=north west, accent] at (0,4.62)
  {phone axes: $z$ out of\\ the screen; swing in the\\ screen plane: $\omega_z=\dot\theta$};
\node[font=\footnotesize, anchor=north west, align=left] at (0,3.62) {phyphox: Gyroscope\\ (rotation rate)};
\begin{scope}[shift={(0.15,1.85)}]
  \draw[->, thin line] (0,0) -- (3.2,0) node[right, font=\scriptsize] {$t$};
  \draw[->, thin line] (0,-0.6) -- (0,0.7) node[right, font=\scriptsize] {$\omega_z$};
  \draw[accent, line width=0.6pt] (0,0.02) -- (0.6,0.02)
    -- plot[domain=0.6:3.05, samples=240] (\x,{-0.55*exp(-0.9*(\x-0.6))*sin(deg(13*(\x-0.6)))});
  \node[font=\scriptsize, note, anchor=south] at (0.3,0.05) {hold};
  \node[font=\scriptsize, note, anchor=west] at (0.75,0.5) {release, free decay};
\end{scope}
\node[caption, anchor=north west] at (0,1.1) {\reltxt};
\end{scope}

% ======================= (b) the measurement problem =======================
\begin{scope}[shift={(9.2,0)}]
\node[panel title] at (0,7.75) {(b) The problem: measure $\zeta$ and $\omega_n$};
\tikzset{cb/.style={block, minimum width=1.3cm, minimum height=0.95cm, inner sep=2pt, font=\footnotesize}}
\node[plant, cb] (pend) at (0.8,6.15) {swing\\ $\theta(0)=\theta_0$};
\node[font=\scriptsize, anchor=north west] at (0.05,5.6) {$\ddot\theta+2\zeta\omega_n\dot\theta+\omega_n^2\theta=0$};
\node[sum] (s) at (2.15,6.15) {};
\draw[sig] (pend) -- node[above] {$\dot\theta$} (s);
\draw[distsig] (2.15,7.05) node[left, font=\footnotesize] {bias $b$, noise $n$} -- (s);
\node[sensor, cb] (gyro) at (3.45,6.15) {gyro $z$\\ quantized};
\draw[sig] (s) -- (gyro);
\node[sensor, cb] (csv) at (5.3,6.15) {phyphox\\ CSV};
\draw[sig] (gyro) -- node[above] {$\omega_z$} (csv);
\node[block, cb] (load) at (5.3,4.35) {pick axis,\\ resample};
\draw[sig] (csv) -- node[right, font=\scriptsize] {$t$, $\omega_{x,y,z}$} (load);
\node[block, cb, minimum width=1.9cm] (int) at (2.95,4.35) {$-\hat b$ (still 3\,s),\\ $\int\!dt$};
\draw[sig] (load) -- (int);
\draw[thick line] (int) -- (1.3,4.35);
\node[font=\footnotesize, anchor=south] at (1.6,4.35) {$\theta(t)$};
\node[ctrl, cb, minimum width=2.3cm] (ld) at (1.3,2.95) {log decrement\\ (peaks)};
\node[ctrl, cb, minimum width=2.3cm] (fit) at (4.0,2.95) {2nd-order fit\\ (free decay)};
\draw[sig] (1.3,4.35) -- (ld);
\draw[sig] (1.3,3.85) -| (fit);
\node[font=\footnotesize] (est) at (6.1,2.95) {$\hat\zeta,\ \hat\omega_n$};
\draw[sig] (fit) -- (est);
\draw[sig] (ld.south) -- ++(0,-0.2) -| (est.south);
% s-plane inset
\begin{scope}[shift={(7.75,5.0)}]
  \draw[->, thin line] (-1.0,0) -- (0.3,0) node[right, font=\scriptsize] {Re};
  \draw[->, thin line] (0,-0.95) -- (0,1.0) node[above, font=\scriptsize] {Im};
  \foreach \y in {0.7,-0.7} \draw[thick line, accent] (-0.3,\y) +(-0.07,-0.07) -- +(0.07,0.07) +(-0.07,0.07) -- +(0.07,-0.07);
  \node[font=\scriptsize, anchor=north, align=center] at (-0.35,-1.0) {poles\\ $-\zeta\omega_n\pm j\omega_d$};
\end{scope}
% goal box
\node[draw=accent, rounded corners=2pt, align=left, font=\footnotesize, anchor=north west, inner sep=4pt,
  text width=8.1cm] at (0,2.05) {\textbf{Estimate} $\zeta,\omega_n$ two ways; compare with $\sqrt{g/d}$, $\sqrt{g/L_\text{eq}}$;\\
  check $T_s\approx 4/(\zeta\omega_n)$; period and damping vs amplitude.\\[2pt]
  \datatxt};
\end{scope}
\end{tikzpicture}
\end{document}
